\section{Introduction}
\label{sec:intro}

Text encoders that produce dense embeddings have become a foundational primitive in modern machine learning, yet the question of how an embedding should \emph{act on} a structured cognitive state has received less attention. We address one instance of this problem: coupling a 384-dimensional French sentence embedding to a Wasserstein gradient flow over a four-species psycholinguistic simplex inspired by the Levelt-Baddeley working-memory model.

Two existing free-energy functionals in our codebase, \texttt{T1FreeEnergy} and \texttt{T2FreeEnergy}, are query-agnostic. Their dynamics depend only on the current state, the species coupling matrix, and a fixed prior. Using either as the oracle for a text-conditioned bridge surrogate produces (state\_pre, state\_post) pairs whose post-state is independent of the query, defeating the purpose of training a text-aware downstream model.

We propose \textsc{QueryConditionedF}, a third free-energy subclass whose accuracy term is parameterized by a JEPA-style decoder \( g_{\mathrm{JEPA}}: \Delta^{128} \to \mathbb{R}^{384} \) that reconstructs the query embedding from the simplex state. The complexity term is the standard per-species KL to a Levelt-Baddeley prior, and a lightweight species-coupling term carries over the inter-code interactions of \texttt{T2FreeEnergy}. The decoder is pre-trained offline against heuristic labels obtained from a deterministic French NLP pipeline (SpaCy-fr, phonemizer, Lexique.org) --- sidestepping the need for ground-truth species annotations.

Our contributions are: (i) a principled formulation of \textsc{QueryConditionedF} with analytical gradients for the complexity and coupling terms and JAX autodiff for the JEPA term; (ii) a reproducible heuristic French labeler producing 32-simplex targets per species; (iii) an empirical ablation of three encoder architectures over a 50k hybrid French corpus, demonstrating that text conditioning shifts the post-state distribution in a predictable, measurable way. A companion full paper (in preparation) extends this design to joint self-supervised training of the encoder and decoder via differentiable JKO.
